% Part: sets-functions-relations
% Chapter: functions
% Section: composition

\documentclass[../../../include/open-logic-section]{subfiles}

\begin{document}

\olfileid{sfr}{fun}{cmp}
\olsection{函数的复合}

\begin{explain}
\oliflabeldef{sfr:fun:inv:sec}{我们在\olref[inv]{sec}节看到，双射~$f$ 的逆~$f^{-1}$ 本身也是一个函数。函数的另一种运算是复合：}{}我们可以通过复合两个函数~$f$ 和~$g$ 来定义新函数，即先应用~$f$ 再应用~$g$。当然，这只有在值域与定义域匹配时才有可能，也就是说，$f$ 的值域必须是~$g$ 的定义域的一个子集。\oliflabeldef{sfr:rel:ops:sec}{函数的这一运算，类似于\olref[rel][ops]{sec}节中关系的相对积运算。}{}

图示可能有助于解释复合的概念。在\olref{fig:composition}中，我们描绘了两个函数~$f \colon A \to B$ 和~$g \colon B \to C$ 以及它们的复合~$(\comp{f}{g})$。函数$(\comp{f}{g}) \colon A \to C$将~$A$ 中的每个元素与~$C$ 中的一个元素配对。我们如下规定~$A$ 中的元素与~$C$ 中的哪个元素配对：给定输入~$x \in A$，首先对~$x$ 应用函数~$f$，它将输出某个~$y \in B$（记作~$f(x) = y$）；然后对~$y$ 应用函数~$g$，它将输出某个$g(f(x)) = g(y) = z \in C$。
\begin{figure}
  \olcachedasset[2\olphotowidth]{composition}
  \caption{函数~$f$ 与~$g$ 的复合~$g \circ f$。}
  \ollabel{fig:composition}
\end{figure}
\end{explain}

\begin{defn}[复合]
设~$f\colon A \to B$ 和~$g\colon B \to C$ 为函数。$f$ 与~$g$ 的\emph{复合}是$\comp{f}{g} \colon A \to C$，其中$(\comp{f}{g})(x) = g(f(x))$。
\end{defn}

\begin{ex}
考虑函数~$f(x) = x + 1$ 和~$g(x) = 2x$。由于$(\comp{f}{g})(x) = g(f(x))$，对每个输入~$x$，需要先取其后继，再将结果乘以二。因此，它们的复合由 $(\comp{f}{g})(x) = 2(x+1)$ 给出。
\end{ex}

\begin{prob}
证明：若~$f \colon A \to B$ 和~$g \colon B \to C$ 都是单射的，则 $\comp{f}{g}\colon A \to C$ 是单射的。
\end{prob}

\begin{prob}
证明：若~$f \colon A \to B$ 和~$g \colon B \to C$ 都是满射的，则 $\comp{f}{g}\colon A \to C$ 是满射的。
\end{prob}

\begin{prob}
设~$f \colon A \to B$ 和~$g \colon B \to C$。证明：$\comp{f}{g}$ 的图象是~$R_f \mid R_g$。
\end{prob}

\end{document}
